\documentclass[11pt, a4paper]{article}

\usepackage[a4paper, top=2.5cm, bottom=2.5cm, left=2cm, right=2cm]{geometry}
\usepackage{amsmath}
\usepackage{amssymb}
\usepackage[colorlinks=true, linkcolor=blue, urlcolor=blue, citecolor=blue]{hyperref}

\title{Theory of Everything - Volume 34 \\ \Large Quantum Computing \& P vs NP}
\author{Omega Protocol}
\date{\today}

\begin{document}

\maketitle

\section*{Layman's Summary}
Standard computers calculate one thing at a time. Quantum computers calculate everything at once. Volume 34 explains how quantum computers tap directly into the fundamental mathematical fabric of the universe, bypassing the limits of normal physics to perform impossible calculations.

\section{Quantum Supremacy (Axiom 37)}
We establish that the computational capacity of the universe doesn't scale normally; it scales exponentially with the size of its "Hilbert space." A quantum computer isn't just a fast calculator; it is a machine that borrows the massive, parallel processing power of the universe's underlying quantum algebra to achieve "Quantum Supremacy."

\section{Shor's Algorithm (Theorem)}
Modern internet security relies on the fact that multiplying two huge numbers is easy, but breaking them back apart (factoring) is incredibly hard. We prove that Shor's Algorithm—which breaks this security instantly—works because it exploits the hidden, periodic "modular flow" of the universe. It finds the hidden rhythms of numbers by vibrating the quantum fabric.

\section{Topological Error Correction (Theorem)}
Quantum computers are notoriously fragile; a slight change in temperature destroys their calculations. We prove that the only way to build a stable quantum computer is through "Topological Error Correction." Instead of storing data in single, fragile atoms, information must be hidden in the global, unbreakable geometry of the quantum network (like tying a knot).

\end{document}
